% !TEX program = xelatex
% Polyp segmentation report (English). Compile: latexmk -xelatex Polyp-Segmentation_Report-EN_v1.tex
% Red text (\ph, \tbd, \figph) marks content still to be written or measured.
\documentclass[10pt,twoside]{article}
\usepackage{report_style}
\usepackage{polyglossia}
\setdefaultlanguage{english}
\gappto\captionsenglish{\renewcommand{\figurename}{Fig.}}

\begin{document}

\makereporttitle
  {Colorectal Polyp Segmentation: Comparing CNN and Transformer Models and Improving Out-of-Domain Generalization}
  {Nguyen An Tran Dien\authorand{and}\ph{[Team member]}}
  {University of Information Technology,\\ Ho Chi Minh City, Vietnam.}

\begin{reportabstract}[Abstract]
\abslabel{Purpose} Polyp segmentation models reach a Dice score close to 0.90 on
images from the centers they were trained on, but lose much of it on images from
other centers: PraNet scores 0.898 on Kvasir-SEG and 0.628 on ETIS. Most published
comparisons also rely on Dice alone. This study compares CNN and Transformer
segmentation models under one training pipeline and tests whether low-cost changes
narrow the out-of-domain gap.
\abslabel{Methods} All models are trained on the PraNet split of Kvasir-SEG and
CVC-ClinicDB (1,450 images) and tested on five sets, three of which come from
centers unseen during training. We reproduce U-Net, MultiResUNet, CE-Net, and
TransUNet \ph{[and Swin-Unet]} with the same data, augmentation, loss, and schedule,
then add a boundary-aware loss and colour-based domain augmentation. Models are
compared with Dice, IoU, recall, and HD95, and differences are tested with the
Wilcoxon signed-rank test on per-image scores.
\abslabel{Results} \ph{[Placeholder] Dice of each model on the in-domain and
out-of-domain sets; size of the generalization gap; whether Transformers beat CNNs
significantly; effect of each improvement in the ablation; inference speed.}
\abslabel{Conclusion} \ph{[Placeholder] One or two sentences on the main finding,
followed by the main limitation (small training set, public benchmarks only, few
runs per configuration).}
\keywords[Keywords]{polyp segmentation, colonoscopy, U-Net, Vision Transformer,
domain generalization, boundary loss, statistical testing}
\end{reportabstract}

% =====================================================================
\section{Introduction}
Most colorectal cancers develop from polyps in the colon, and finding and removing
polyps during colonoscopy lowers the risk of cancer \ph{[citation]}. Polyps vary in
size, shape, and colour, and many look similar to the surrounding mucosa, so some
are missed. Automatic segmentation outlines the polyp pixel by pixel and can support
the endoscopist during the procedure. Encoder--decoder networks derived from
U-Net~\cite{ronneberger2015unet} are the standard approach, and recent models extend
or replace the convolutional encoder with
Transformers~\cite{chen2021transunet,cao2022swinunet}.

Scores on the usual benchmarks are high, but they drop sharply on data from other
centers. PraNet reaches a mean Dice of 0.898 on Kvasir-SEG and 0.899 on
CVC-ClinicDB, the two sources of its training data, but 0.709 on CVC-ColonDB and
0.628 on ETIS~\cite{fan2020pranet}. Evaluation practice adds a second problem. An
audit of 27 polyp segmentation papers found that 26 report no significance test,
25 omit HD95, and 25 omit recall; on the six unseen centers of PolypGen, all
re-evaluated models fell to a Dice of 0.71--0.76~\cite{urooj2026audit}. Small Dice
differences measured on in-domain sets therefore say little about which model works
better on new data.

This report does not attempt to set a new state of the art. It studies a narrower
question: under one controlled pipeline, how do the CNN and Transformer models of
the course reading list~\cite{ronneberger2015unet,chen2021transunet,cao2022swinunet,ibtehaz2020multiresunet,gu2019cenet,isensee2018nnunet}
compare on in-domain and out-of-domain data, and can low-cost changes reduce the gap
between the two?

The report addresses three research questions:
\begin{itemize}
  \item \textbf{RQ1.} With about 1,450 training images, do Transformer-based models
  (TransUNet, Swin-Unet) outperform CNN-based models (MultiResUNet, CE-Net) by a
  statistically significant margin?
  \item \textbf{RQ2.} Which model keeps its performance best on the out-of-domain
  sets CVC-ColonDB, ETIS, and CVC-300?
  \item \textbf{RQ3.} How much do a boundary-aware loss and colour-based domain
  augmentation reduce the generalization gap, and which component contributes most?
\end{itemize}

The report is organized as follows. Section~\ref{sec:related} reviews encoder--decoder
CNNs, Transformer-based segmentation, and polyp benchmarks. Section~\ref{sec:problem}
defines the task and the data. Section~\ref{sec:pipeline} describes the experimental
pipeline and the models, followed by the loss functions and the proposed improvements
in Section~\ref{sec:methods}. Sections~\ref{sec:setup} and~\ref{sec:results} present
the experimental setup and results. Section~\ref{sec:limits} states the reliability
constraints, and Section~\ref{sec:conclusion} concludes the study.

% =====================================================================
\section{Related Works}\label{sec:related}
\subsection{Encoder--decoder CNNs for medical image segmentation}
U-Net combines a contracting encoder with an expanding decoder and passes
high-resolution features through skip connections, which lets it learn from few
annotated images~\cite{ronneberger2015unet}. MultiResUNet replaces the two
convolutions of each stage with a MultiRes block, a chain of $3\times3$ convolutions
whose outputs are concatenated to cover several receptive fields, and adds residual
``Res paths'' to the skip connections to reduce the semantic gap between encoder and
decoder features; it was evaluated on five biomedical datasets, including
CVC-ClinicDB and ISIC~2018~\cite{ibtehaz2020multiresunet}. CE-Net places a dense
atrous convolution (DAC) block and a residual multi-kernel pooling (RMP) block
between a pretrained ResNet encoder and the decoder, so that multi-scale context is
captured without further downsampling~\cite{gu2019cenet}. nnU-Net proposes no new
architecture: it configures preprocessing, network topology, training, and
post-processing from a ``fingerprint'' of the dataset and was evaluated on the ten
tasks of the Medical Segmentation Decathlon~\cite{isensee2018nnunet}.

\subsection{Transformers for medical image segmentation}
Convolutions capture local patterns, and their limited receptive field makes
long-range context hard to model. TransUNet passes CNN feature maps to a Vision
Transformer encoder and decodes them with a cascaded upsampler and U-Net skip
connections~\cite{chen2021transunet}. Swin-Unet builds the whole U-shaped network
from Swin Transformer blocks, whose shifted windows reduce the cost of
self-attention, and upsamples with patch-expanding layers~\cite{cao2022swinunet}.
Both models were evaluated on the Synapse multi-organ CT and ACDC cardiac MRI
datasets rather than on colonoscopy images.

\subsection{Polyp segmentation benchmarks and evaluation practice}
PraNet introduced the protocol that most polyp papers follow: training on 900
Kvasir-SEG and 550 CVC-ClinicDB images and testing on five
sets~\cite{fan2020pranet}. PolypGen adds images from six centers to test
generalization~\cite{ali2023polypgen}. The audit in~\cite{urooj2026audit} found at
least five incompatible split protocols in use, transcription errors in survey
tables, and differences of up to about 16~px in HD95 between metric libraries. It
proposes a reporting checklist: HD95 and recall as primary metrics, an explicit split
protocol, an external test set, and significance tests over per-image scores. Metrics
Reloaded gives general guidance on choosing segmentation
metrics~\cite{maierhein2024metrics}. This report follows these recommendations.

% =====================================================================
\section{Problem Formulation and Data}\label{sec:problem}
\subsection{Task formulation}
An image $x\in\mathbb{R}^{H\times W\times 3}$ has a binary mask
$y\in\{0,1\}^{H\times W}$ that marks the polyp pixels. A model $f_\theta$ outputs a
probability map, and the predicted mask is obtained by thresholding:
\begin{equation}
  p = f_\theta(x)\in[0,1]^{H\times W},\qquad \hat{y} = \mathbb{1}[\,p \ge 0.5\,].
\end{equation}
Training images come from a source domain $\mathcal{D}_S$ (Kvasir-SEG and
CVC-ClinicDB). Test images come from $\mathcal{D}_S$ and from $K$ target domains
$\mathcal{D}_T^1,\dots,\mathcal{D}_T^K$ that are never seen during training
(CVC-ColonDB, ETIS, and CVC-300). For a per-image score $s$ (Dice in the main
analysis), the generalization gap of a model is
\begin{equation}\label{eq:gap}
  \Delta = \bar{s}\bigl(\mathcal{D}_S^{\text{test}}\bigr)
  - \frac{1}{K}\sum_{k=1}^{K}\bar{s}\bigl(\mathcal{D}_T^{k}\bigr),
\end{equation}
where $\bar{s}$ is the mean score over the images of a set. A model generalizes
better when it scores high on the target domains and has a small $\Delta$.

\subsection{Data sources}
Table~\ref{tab:data} lists the five public datasets. Kvasir-SEG and CVC-ClinicDB
provide the training images; CVC-ColonDB, ETIS, and CVC-300 are used only for
testing. All five are used in the form released with PraNet~\cite{fan2020pranet}.

\begin{table}[t]
\centering
\caption{Datasets and their roles.}\label{tab:data}
\begin{tabular}{@{}llll@{}}
\toprule
Dataset & Images & Resolution & Role \\
\midrule
Kvasir-SEG~\cite{jha2020kvasir} & 1,000 & varied & Training; in-domain test \\
CVC-ClinicDB~\cite{bernal2015clinicdb} & 612 & $384\times288$ & Training; in-domain test \\
CVC-ColonDB~\cite{bernal2012colondb} & 380 & $574\times500$ & Out-of-domain test \\
ETIS-LaribPolypDB~\cite{silva2014etis} & 196 & $1225\times966$ & Out-of-domain test \\
CVC-300 (EndoScene)~\cite{vazquez2017endoscene} & 60 & \ph{[verify]} & Out-of-domain test \\
\bottomrule
\end{tabular}
\end{table}

\subsection{Data split}
Table~\ref{tab:split} gives the number of images per split. Training and test images
follow the fixed PraNet split~\cite{fan2020pranet}. Ten percent of the training
images of each source dataset are held out at random (seed 42) for validation, so
that checkpoints are selected without touching any test set. Because the three CVC
datasets were collected by related groups, near-duplicate images across the training
and test sets are checked with perceptual hashing before training, as
in~\cite{urooj2026audit}. Figure~\ref{fig:samples} shows sample images.

\begin{table}[t]
\centering
\caption{Number of images per split.}\label{tab:split}
\begin{tabular}{@{}lrrrrrr@{}}
\toprule
Split & Kvasir & ClinicDB & ColonDB & ETIS & CVC-300 & Total \\
\midrule
Training   & 810 & 495 & -- & -- & -- & 1,305 \\
Validation &  90 &  55 & -- & -- & -- &   145 \\
Test       & 100 &  62 & 380 & 196 & 60 & 798 \\
\bottomrule
\end{tabular}
\end{table}

\begin{figure}[t]
\centering
\figph{[Figure placeholder] One image and its mask from each of the five datasets,
showing differences in colour, lighting, and resolution. Generated by
scripts/plot\_dataset\_examples.py.}
\caption{Sample images and ground-truth masks from each dataset.}\label{fig:samples}
\end{figure}

% =====================================================================
\section{Experimental Pipeline}\label{sec:pipeline}
Figure~\ref{fig:pipeline} shows the pipeline, and Table~\ref{tab:models} lists the
models. Every model goes through the same stages with the same data, augmentation,
loss, and schedule; only the architecture changes, so differences in the results can
be attributed to the architecture rather than to the training recipe. The shaded
stages are where the improvements of Section~\ref{sec:methods} are applied.

\begin{figure}[t]
\centering
\begin{tikzpicture}[
  every node/.append style={execute at begin node={\hyphenpenalty=10000}},
  box/.style={draw, rounded corners=2pt, align=center, font=\scriptsize,
    text width=72pt, minimum height=40pt, inner sep=2pt},
  hl/.style={box, fill=black!8},
  arr/.style={-{Stealth[length=3.5pt]}, semithick}]
  \node[box] (a) at (40pt,0)    {\textbf{1 Data}\\ PraNet split\\ 5 datasets};
  \node[box] (b) at (136pt,0)   {\textbf{2 Preprocessing}\\ resize, ImageNet\\ normalization};
  \node[hl]  (c) at (232pt,0)   {\textbf{3 Augmentation}\\ flips, rotation,\\ scale, colour (C3)};
  \node[box] (d) at (328pt,0)   {\textbf{4 Model}\\ U-Net, CE-Net,\\ MultiResUNet,\\ TransUNet, Swin-Unet};
  \node[hl]  (e) at (328pt,-56pt) {\textbf{5 Training}\\ BCE + Dice,\\ boundary loss (C2)};
  \node[hl]  (f) at (232pt,-56pt) {\textbf{6 Inference}\\ threshold 0.5,\\ test-time aug. (C3)};
  \node[box] (g) at (136pt,-56pt) {\textbf{7 Evaluation}\\ Dice, IoU,\\ recall, HD95};
  \node[box] (h) at (40pt,-56pt)  {\textbf{8 Statistics, demo}\\ Wilcoxon and\\ Friedman tests;\\ Gradio demo};
  \draw[arr] (a) -- (b); \draw[arr] (b) -- (c); \draw[arr] (c) -- (d);
  \draw[arr] (d) -- (e); \draw[arr] (e) -- (f); \draw[arr] (f) -- (g); \draw[arr] (g) -- (h);
  \node[draw, dashed, rounded corners=2pt, font=\scriptsize, text width=356pt,
    align=center, inner sep=4pt] at (184pt,-102pt)
    {Same data, schedule, and seed for every model \,$\bullet$\, checkpoint chosen on
     validation \,$\bullet$\, per-image scores logged for every test set};
\end{tikzpicture}
\caption{Experimental pipeline. Shaded stages are modified by the proposed
improvements (C2: boundary-aware loss; C3: domain augmentation and test-time
augmentation).}\label{fig:pipeline}
\end{figure}

\begin{table}[t]
\centering
\caption{Models compared in this study (``\#'' refers to the course reading list).}\label{tab:models}
\begin{tabularx}{\linewidth}{@{}l c >{\raggedright\arraybackslash}X >{\raggedright\arraybackslash}X l@{}}
\toprule
Model & \# & Key idea & Problem addressed & Priority \\
\midrule
U-Net~\cite{ronneberger2015unet} & -- & Encoder--decoder with skip connections & Reference baseline & Required \\
MultiResUNet~\cite{ibtehaz2020multiresunet} & 59 & MultiRes blocks; Res paths & Polyps of varying size; encoder--decoder semantic gap & Required \\
CE-Net~\cite{gu2019cenet} & 60 & DAC and RMP blocks at the bottleneck & Spatial detail lost by repeated pooling & Required \\
TransUNet~\cite{chen2021transunet} & 57 & CNN features, ViT encoder, U-Net decoder & Limited receptive field of CNNs & Required$^{*}$ \\
Swin-Unet~\cite{cao2022swinunet} & 56 & U-shaped network of Swin Transformer blocks & Global context at lower cost than ViT & Optional$^{*}$ \\
nnU-Net~\cite{isensee2018nnunet} & 58 & Self-configuring U-Net pipeline & Strong baseline without manual tuning & Optional \\
\bottomrule
\multicolumn{5}{@{}l}{$^{*}$At least one of the two Transformer models is required.}
\end{tabularx}
\end{table}

TransUNet and Swin-Unet expect $224\times224$ inputs by default, while the CNNs follow
PraNet and use $352\times352$. The input size of each model is reported in
Appendix~\ref{app:settings}, and predictions are resized to the original resolution
before evaluation, so all models are scored on the same masks. Each run writes its
configuration, seed, and per-image scores to a shared log. The best model is served
through a Gradio~\cite{abid2019gradio} demo that overlays the predicted mask on an
uploaded image.

% =====================================================================
\section{Loss Functions and Proposed Improvements}\label{sec:methods}
\subsection{Segmentation loss}
All models are trained with the sum of binary cross-entropy and a soft Dice loss:
\begin{equation}\label{eq:loss}
  \mathcal{L}_{\text{seg}} = \mathcal{L}_{\text{BCE}}(p,y)
  + 1 - \frac{2\sum_i p_i y_i + \epsilon}{\sum_i p_i + \sum_i y_i + \epsilon},
\end{equation}
where $i$ runs over pixels and $\epsilon$ avoids division by zero. Cross-entropy gives
a gradient for every pixel, and the Dice term counteracts the imbalance between the
small polyp region and the large background.

\subsection{Boundary-aware loss (C2)}
Dice and cross-entropy count pixels, so a prediction that covers the polyp but places
its edge a few pixels off is barely penalized, whereas HD95 is sensitive to such
errors. Following Kervadec et al.~\cite{kervadec2021boundary}, we add a term weighted
by the signed distance map $\phi_G$ of the ground-truth boundary, negative inside the
polyp and positive outside:
\begin{equation}\label{eq:boundary}
  \mathcal{L}_{B} = \frac{1}{|\Omega|}\sum_{q\in\Omega}\phi_G(q)\,p(q),
  \qquad
  \mathcal{L} = \mathcal{L}_{\text{seg}} + \lambda\,\mathcal{L}_{B}.
\end{equation}
Probability placed far outside the true region costs more than probability near its
edge. As recommended in~\cite{kervadec2021boundary}, $\lambda$ starts small and is
increased during training, up to \ph{[value]}. An auxiliary head that predicts the
boundary map is a fallback if this term makes training unstable.

\subsection{Domain augmentation and test-time augmentation (C3)}
We assume that a large part of the difference between centers comes from the colour,
illumination, and resolution of the endoscopes. The baseline augmentation (flips,
rotation, scaling, mild brightness changes) is therefore extended with stronger
photometric changes: brightness and contrast, hue and saturation shifts, gamma, blur,
and JPEG compression, implemented with Albumentations~\cite{buslaev2020albumentations}.
At inference time, predictions are averaged over horizontal and vertical flips. The
effect is measured by the change in the gap $\Delta$ of Eq.~\eqref{eq:gap}.

\subsection{Stronger pretrained encoder (C1, optional)}
If the baseline comparison is complete by week 4 of the project, the encoder of the
best model is replaced by a PVTv2~\cite{wang2022pvtv2} or ConvNeXt~\cite{liu2022convnext}
backbone pretrained on ImageNet, while the decoder keeps the Res paths of
MultiResUNet. The aim is to combine global context with multi-scale detail.

% =====================================================================
\section{Experimental Setup}\label{sec:setup}
\begin{table}[t]
\centering
\caption{Main experimental configuration.}\label{tab:config}
\begin{tabularx}{\linewidth}{@{}l >{\raggedright\arraybackslash}X@{}}
\toprule
Component & Configuration \\
\midrule
Training data & 1,305 images (PraNet training split minus 10\% for validation) \\
Input size & $352\times352$ for CNNs; $224\times224$ for TransUNet and Swin-Unet \ph{[confirm]} \\
Normalization & ImageNet mean and standard deviation \\
Augmentation & Flips, rotation, random scaling, mild colour changes; stronger photometric set for C3 \\
Optimizer & AdamW~\cite{loshchilov2019adamw}, learning rate $10^{-4}$, weight decay \ph{[value]} \\
Batch size, epochs & \ph{8--16}, 100 epochs, mixed precision \\
Loss & BCE + Dice (Eq.~\ref{eq:loss}); boundary term for C2 (Eq.~\ref{eq:boundary}) \\
Checkpoint selection & Highest validation Dice; test sets never used for selection \\
Framework, hardware & PyTorch~\cite{paszke2019pytorch} \ph{[version]}; Google Colab or Kaggle GPU \ph{[model]} \\
Seeds & 42 \ph{[and two more if GPU time allows]} \\
\bottomrule
\end{tabularx}
\end{table}

Table~\ref{tab:config} summarizes the configuration. For a predicted mask $\hat{y}$ and
a ground-truth mask $y$, the overlap and sensitivity metrics are
\begin{equation}
  \text{Dice} = \frac{2|\hat{y}\cap y|}{|\hat{y}|+|y|},\qquad
  \text{IoU} = \frac{|\hat{y}\cap y|}{|\hat{y}\cup y|},\qquad
  \text{Recall} = \frac{|\hat{y}\cap y|}{|y|}.
\end{equation}
Recall is reported because a missed polyp is more harmful than a false alarm. Boundary
accuracy is measured with the 95th percentile of the symmetric Hausdorff distance,
\begin{equation}
  \text{HD95}(\hat{y},y) = \max\Bigl\{ P_{95}\bigl\{d(a,\partial y) : a\in\partial\hat{y}\bigr\},\;
  P_{95}\bigl\{d(b,\partial\hat{y}) : b\in\partial y\bigr\}\Bigr\},
\end{equation}
where $\partial$ denotes the boundary of a mask and $d$ the Euclidean distance in pixels
at the original resolution. Since HD95 values differ between
libraries~\cite{urooj2026audit}, a single library (\ph{[name, version]}) is used for
all models, and an empty prediction receives the image diagonal as its HD95. All
metrics are computed per image and averaged per test set (mDice, mIoU).

Two models are compared with the Wilcoxon signed-rank test~\cite{wilcoxon1945} on
paired per-image Dice scores; a difference is called significant only when $p<0.05$.
To compare all models at once, the Friedman test~\cite{friedman1937} is followed by a
post-hoc test, as recommended by Dem\v{s}ar~\cite{demsar2006}. Means on the
out-of-domain sets are reported with 95\% bootstrap confidence intervals. Efficiency is
reported as the number of parameters, FLOPs, and frames per second at batch size 1.

% =====================================================================
\section{Results and Discussion}\label{sec:results}
\subsection{CNN versus Transformer models (RQ1)}
\ph{[Placeholder] Compare mDice of all models on the five test sets (Table~\ref{tab:main}).
State whether the best Transformer is significantly better than the best CNN, and
discuss the role of ImageNet pretraining.}

\begin{table}[t]
\centering
\caption{mDice on the five test sets. PraNet values are those reported
in~\cite{fan2020pranet} and are given for reference only.}\label{tab:main}
\begin{tabular}{@{}lcccccc@{}}
\toprule
 & \multicolumn{2}{c}{In-domain} & \multicolumn{3}{c}{Out-of-domain} & \\
\cmidrule(lr){2-3}\cmidrule(lr){4-6}
Model & Kvasir & ClinicDB & ColonDB & ETIS & CVC-300 & $p$ vs.\ CNN \\
\midrule
PraNet~\cite{fan2020pranet} & 0.898 & 0.899 & 0.709 & 0.628 & 0.871 & -- \\
\midrule
U-Net        & \tbd & \tbd & \tbd & \tbd & \tbd & \tbd \\
MultiResUNet & \tbd & \tbd & \tbd & \tbd & \tbd & \tbd \\
CE-Net       & \tbd & \tbd & \tbd & \tbd & \tbd & \tbd \\
TransUNet    & \tbd & \tbd & \tbd & \tbd & \tbd & \tbd \\
Swin-Unet    & \tbd & \tbd & \tbd & \tbd & \tbd & \tbd \\
\bottomrule
\end{tabular}
\end{table}

\subsection{Out-of-domain generalization (RQ2)}
\ph{[Placeholder] Report the gap $\Delta$, recall, and HD95 on the out-of-domain sets
(Table~\ref{tab:ood}). Identify which model degrades least and whether the ranking
changes between Dice and HD95 (Figure~\ref{fig:gap}).}

\begin{table}[t]
\centering
\caption{Generalization of each model. In-domain: Kvasir and ClinicDB; out-of-domain:
ColonDB, ETIS, and CVC-300 (mean over the three sets, 95\% CI in brackets).}\label{tab:ood}
\begin{tabular}{@{}lccccc@{}}
\toprule
Model & mDice in & mDice out & $\Delta$ & Recall out & HD95 out (px) \\
\midrule
U-Net        & \tbd & \tbd & \tbd & \tbd & \tbd \\
MultiResUNet & \tbd & \tbd & \tbd & \tbd & \tbd \\
CE-Net       & \tbd & \tbd & \tbd & \tbd & \tbd \\
TransUNet    & \tbd & \tbd & \tbd & \tbd & \tbd \\
Swin-Unet    & \tbd & \tbd & \tbd & \tbd & \tbd \\
\bottomrule
\end{tabular}
\end{table}

\begin{figure}[t]
\centering
\figph{[Figure placeholder] Left: mDice in-domain vs.\ out-of-domain for each model
(paired dots). Right: distribution of per-image Dice on ETIS.}
\caption{Generalization gap of each model.}\label{fig:gap}
\end{figure}

\subsection{Ablation of the proposed improvements (RQ3)}
\ph{[Placeholder] Starting from the best baseline after week 4, add each component and
report its effect (Table~\ref{tab:ablation}). Report an improvement that is not
significant as a negative result.}

\begin{table}[t]
\centering
\caption{Ablation on the best baseline model. $p$-values compare each row with the
baseline on per-image out-of-domain Dice.}\label{tab:ablation}
\begin{tabular}{@{}lccccc@{}}
\toprule
Configuration & mDice in & mDice out & $\Delta$ & HD95 out (px) & $p$ \\
\midrule
Best baseline \ph{[name]}         & \tbd & \tbd & \tbd & \tbd & -- \\
+ C3 (domain aug.\ + TTA)         & \tbd & \tbd & \tbd & \tbd & \tbd \\
+ C2 (boundary loss)              & \tbd & \tbd & \tbd & \tbd & \tbd \\
+ C2 + C3                         & \tbd & \tbd & \tbd & \tbd & \tbd \\
C1 + C2 + C3 (optional)           & \tbd & \tbd & \tbd & \tbd & \tbd \\
\bottomrule
\end{tabular}
\end{table}

\subsection{Error analysis and efficiency}
\ph{[Placeholder] Show the lowest-Dice images of each test set (Figure~\ref{fig:errors}),
group results by polyp size, and compare model cost (Table~\ref{tab:cost}).}

\begin{table}[t]
\centering
\caption{Model size and inference speed (batch size 1, \ph{[GPU]}).}\label{tab:cost}
\begin{tabular}{@{}lccc@{}}
\toprule
Model & Parameters (M) & FLOPs (G) & Frames per second \\
\midrule
U-Net        & \tbd & \tbd & \tbd \\
MultiResUNet & \tbd & \tbd & \tbd \\
CE-Net       & \tbd & \tbd & \tbd \\
TransUNet    & \tbd & \tbd & \tbd \\
Swin-Unet    & \tbd & \tbd & \tbd \\
Proposed     & \tbd & \tbd & \tbd \\
\bottomrule
\end{tabular}
\end{table}

\begin{figure}[t]
\centering
\figph{[Figure placeholder] Image, ground truth, and predictions of each model for the
three lowest-Dice images of ETIS and CVC-ColonDB.}
\caption{Typical failure cases on out-of-domain images.}\label{fig:errors}
\end{figure}

% =====================================================================
\section{Reliability and Limitations}\label{sec:limits}
The evidence is constrained in several ways. The models are trained on 1,305 images
from two sources, and the out-of-domain sets are small (60--380 images), so their means
carry wide confidence intervals; differences smaller than these intervals are not
interpreted. All test images contain at least one polyp, so false alarms on normal
frames are not measured, and still images selected by dataset creators are cleaner than
live video, which contains blur, specular highlights, and motion.

Free GPU time limits the number of runs, so some configurations are trained with a single
seed. Hyperparameters are shared rather than tuned per model, which may favour models
whose original recipe resembles the shared one. The reproduced models are trained on
different data and input sizes than in their original papers, so the scores are compared
only within this pipeline and with PraNet, not with the numbers reported in the original
papers. The study is an academic comparison on public data; it involves no clinical
validation, and the demo is not a medical device.

% =====================================================================
\section{Conclusion and Future Work}\label{sec:conclusion}
\ph{[Placeholder] Summarize the answers to RQ1--RQ3 in one paragraph, with the main
numbers.}

Future work should test the models on PolypGen~\cite{ali2023polypgen}, whose six centers
give a stronger external evaluation than the public benchmark sets. Evaluating on
colonoscopy videos would measure stability across frames and false alarms on frames
without polyps. Lighter models suited to real-time inference, and semi-supervised
training with unlabelled colonoscopy frames, are further extensions.

\runin{Data and Code Availability.} The project code, configurations, data split lists,
per-image scores, and plotting scripts are maintained in the project repository
\ph{[URL]}. The datasets are public and remain subject to their original terms; they are
not redistributed.

% =====================================================================
\startappendix[Appendix]
\section{Model-Specific Settings}\label{app:settings}
\begin{table}[h]
\centering
\caption{Settings that differ between models (IN: ImageNet). All other
settings follow Table~\ref{tab:config}.}\label{tab:settings}
\setlength{\tabcolsep}{4pt}
\begin{tabular}{@{}lccccc@{}}
\toprule
Model & Input size & Pretrained weights & Batch & Epochs & Params (M) \\
\midrule
U-Net        & $352\times352$ & none & \tbd & \tbd & \tbd \\
MultiResUNet & $352\times352$ & none & \tbd & \tbd & \tbd \\
CE-Net       & $352\times352$ & ResNet-34 (IN-1k) & \tbd & \tbd & \tbd \\
TransUNet    & $224\times224$ & R50-ViT-B/16 (IN-21k) & \tbd & \tbd & \tbd \\
Swin-Unet    & $224\times224$ & Swin-T (IN-1k) & \tbd & \tbd & \tbd \\
Proposed     & \tbd & \tbd & \tbd & \tbd & \tbd \\
\bottomrule
\end{tabular}
\end{table}

\clearpage
\section{Reproducibility Artifact Map}\label{app:artifacts}
\begin{table}[h]
\centering
\caption{Primary artifacts in the project repository (planned layout; update to the
final paths).}\label{tab:artifacts}
\begin{tabular}{@{}ll@{}}
\toprule
Purpose & Artifact \\
\midrule
Run instructions        & README.md \\
Data split lists        & data/splits/\{train,val,test\}\_*.txt \\
Duplicate check         & scripts/check\_duplicates.py \\
Experiment configurations & configs/*.yaml \\
Model implementations   & src/models/ \\
Losses (BCE + Dice, boundary) & src/losses.py \\
Training and evaluation & scripts/train.py, scripts/evaluate.py \\
Per-image scores        & outputs/metrics/per\_image\_*.csv \\
Statistical tests       & scripts/stats\_tests.py \\
Report figures          & scripts/plot\_report\_figures.py $\rightarrow$ report/figures/ \\
Demo                    & app/gradio\_app.py \\
\bottomrule
\end{tabular}
\end{table}

\printreferences

\end{document}
